\scanpage{40}
% archive foliation: 39
\uline{Prop} Soient $M,N$ deux $K$-modules, $u$ $(v)$ un endom de $M$ $(N)$, $f:M\longrightarrow N$ un $K$-hom. Considérons $M,N$ comme $\overline G_{aK}$-modules par $\lambda\in\overline G_{aK}(K')$ opérant par $\exp(\lambda u)$ resp. $\exp(\lambda v)$. Pour que $f$ soit un $G_{aK}$-morphisme, il f. et s. que l'on ait $vf=fu$ i.e. $vf-fu=0$

\uline{Cor} Si $u,v$ \deleted{\unsure{c|et}} sont nilpotents, c'est aussi la cond. néc. et suff. pour que $f$ soit un $G_a$-\unsure{morphisme}.

\uline{Dém} Il faut écrire
\[
\sum_{n\geq0}\lambda^n\frac{v^n}{n!}f=\sum_{n\geq0}\lambda^n f\frac{u^n}{n!}
\]
\margin{\unsure{univ\textsuperscript{t}.|univ.} en $\lambda$ nilp.}
i.e.
\[
\sum_{n\geq1}\frac1{n!}(v^nf-fu^n)\lambda^n=0
\]
\unsure{univ.} en $\lambda$ nilp., ce qui équivaut à
\[
v^nf-fu^n=0\qquad\forall n\geq1
\]
Pour $n=1$, c'est la commutation $vf=fu$, qui implique par récurrence $v^nf=fu^n$ $(\forall n\geq1)$.

\uline{Proposition} Soient \deleted{\unsure{u,v}} $u\in N(\operatorname{End}_K(M))$, $v\in N\operatorname{End}_K(N)$, $f\in\operatorname{Hom}_K(M,N)$. Conditions équivalentes :

1) $f$ est un $\overline G_{aK}$-hom (pour \inserted{$\overline G_a$ opérant sur $M,N$ par} $\bar e_u$ et $\bar e_v$)

2) $f$ est un $G_{aK}$-hom ($\ldots$)

3) $fu=vf\deleted{\unsure{0}}$ i.e. $\delta\overset{\mathrm{déf}}=vf-fu=0$

4) $f\circ\exp(u)=\exp(v)\circ f$
